Cette section vise à décrire tout le matériel mis à disposition, 
ainsi que la description de manière séparée 
des trois manipulations qui ont été réalisées.

\subsection{Matériel utilisé}
Le dispositif expérimental repose sur les éléments suivants :
\begin{itemize}
    \item Un banc optique gradué pour la mesure des positions ;
    \item Une source de lumière blanche associée à un objet ;
    \item Un écran quadrillé pour l'observation des images ;
    \item Une lentille mince convergente $L_1$ de focale $f'_1$ inconnue ;
    \item Une lentille divergente $L \prime$ de focale $f$ inconnue ; 
    \item Un miroir plan. 
\end{itemize}

\subsection{Manipulation 1 : Lentilles convergentes}

\subsubsection{Image à l'infini et Autocollimation}
Pour la mesure par image à l'infini :
\begin{enumerate}
    \item Positionner l'écran à l'extrémité du banc optique.
    \item Déplacer la lentille $L_1$ entre l'objet et l'écran jusqu'à obtenir une image nette.
    \item Relever la position de la lentille pour déterminer la distance lentille-écran $f'_1$.
\end{enumerate}

Pour l'autocollimation :
\begin{enumerate}
    \item Placer un miroir plan immédiatement derrière la lentille $L_1$.
    \item Déplacer l'ensemble \{lentille + miroir\} le long du banc optique.
    \item Identifier la position où l'image de l'objet est nette dans le plan de l'objet.
    \item Mesurer la distance entre le centre optique et l'objet pour obtenir $f'_1$.
\end{enumerate}

\subsubsection{Méthode de Silbermann}
\begin{enumerate}
    \item Installer l'objet et l'écran de part et d'autre de la lentille $L_1$.
    \item Déplacer conjointement la lentille et l'écran pour obtenir une image inversée.
    \item Ajuster jusqu'à ce que la taille de l'image soit strictement identique à celle de l'objet.
    \item Relever les distances $p$ et $p'$.
\end{enumerate}

\subsubsection{Méthode de Bessel}
\begin{enumerate}
    \item Fixer l'objet et l'écran à une distance $D$ constante respectant $D~=~50$~cm.
    \item Faire coulisser la lentille $L_1$ entre ces deux éléments fixes.
    \item Noter la position $O_1$ pour laquelle l'image est nette et agrandie.
    \item Noter la position $O_2$ pour laquelle l'image est nette et rétrécie.
    \item Mesurer la distance $d$ séparant ces deux positions de la lentille.
\end{enumerate}

\subsubsection{Méthode des points conjugués}
\begin{enumerate}
    \item Placer l'objet à différentes distances $p$ de la lentille (de $-30$ cm à $-15$ cm).
    \item Chercher, pour chaque valeur, la position de l'écran assurant une image nette.
    \item Relever systématiquement le couple de distances $(p, p')$.
    \item Pour le cas "objet virtuel", utiliser $L_2$ pour faire converger le faisceau derrière $L_1$.
\end{enumerate}

\subsection{Manipulation 2 : Méthode de Badal}
\begin{enumerate}
    \item Former une image à l'infini de l'objet grâce à la lentille $L_1$.
    \item Positionner la lentille auxiliaire $L_2$ pour observer cette image sur l'écran.
    \item Relever la position de référence de l'écran $d_1$.
    \item Intercaler la lentille divergente $L'$ au niveau du foyer objet de $L_2$.
    \item Ajuster l'écran pour retrouver la netteté et noter la nouvelle position $d_2$.
\end{enumerate}